% ===================================================================
%  TABLICA Nr 12 — Dworzec. — Koleje. — Statki.
%  Meme regime que les precedents. Tableau a UNE SEULE COLONNE.
%
%  DEUX NOTES, ET ELLES NE SONT PAS AU MEME ENDROIT. Celle de
%  l'apparat — les horaires different d'un pays a l'autre, on prend
%  ceux du tableau francais — ouvre le tableau ; celle en pied,
%  t12-noto-1, dit que le voyage est decrit selon la frontiere
%  d'avant-guerre, et le fac-simile la place apres le 14. Le
%  neerlandais la laisse la ; cette colonne aussi. Le rang du bloc
%  fait partie du document.
%
%  LE 11 ET LE 18 ONT UN RENVOI HORS DU GRAS : « pedvarmigili » (48)
%  et « sennombra » (88) sont composes en romain, renvoi compris. Ce
%  n'est pas un oubli a reparer — le releve compte les renvois, pas
%  les gras, et une colonne qui rajouterait le gras dirait autre
%  chose que la planche.
%
%  ET DEUX MOTS DEMANDENT DE NE PAS SE FIER AU FRANCAIS :
%  « biblioteko » (14) n'est pas la bibliotheque mais le kiosque a
%  livres du hall, et « trotuaro » (53, 65) n'est pas le trottoir
%  mais le quai. Le polonais dit ksiegarnia et peron.
% ===================================================================

%%K t12-apar-1 apar
\VUsaut{4.0mm}
\VUtitre{15.0pt}{60}{TABLICA Nr 12}
\VUsaut{2.4mm}
\VUpk{11.4pt}{40}{Dworzec. --- Koleje. --- Statki.}
\VUsaut{3.4mm}
\textsc{Uwaga.} --- \textit{Rozkłady godzin używane w każdym kraju są rozmaite, bierzemy
więc za przykład godziny} \VUgras{tablicy francuskiej.}

%%K t12-01-1 p
1. --- Właśnie podróżowałem przez Niemcy. Przed odjazdem kupiłem \VUgras{książkę
rozkładu jazdy} \textsuperscript{(15)}, żeby poznać godzinę odjazdu pociągów.
Stwierdziwszy, że najdogodniejszy pociąg ma odjechać z Paryża o dziewiątej wieczorem i
przybyć do Strasburga o pierwszej trzynaście, przygotowałem swoją podróż, a nazajutrz
kazałem się zawieźć na dworzec.

%%K t12-02-1 p
2. --- Kiedy wszedłem do wielkiej hali odjazdów, za starym wiejskim \VUgras{księdzem}
\textsuperscript{(2)}, który niósł w ręku swoją \VUgras{torbę podróżną}
\textsuperscript{(3)}, wielu \VUgras{podróżnych} \textsuperscript{(1)} już przybyło.

%%K t12-02-2 p
Ponieważ chciałem wysłać \VUgras{depeszę (telegram)} \textsuperscript{(5)}, kazałem
sobie wskazać przez \VUgras{kontrolera} \textsuperscript{(6)} \VUgras{biuro
telegraficzne} \textsuperscript{(4)}.

%%K t12-03-1 p
3. --- Zanim wszedłem do \VUgras{poczekalni} \textsuperscript{(7)}, poszedłem zaopatrzyć
się w \VUgras{bilet} \textsuperscript{(9)} powrotny.

%%K t12-04-1 p
4. --- Nie czekałem długo przy \VUgras{okienku} \textsuperscript{(8)}, bo mało osób tam
było. \VUgras{Żołnierz} \textsuperscript{(10)}, który odjeżdżał na urlop, uprzejmie
ustąpił mi swoją kolej, tak że miałem dość czasu, żeby poprosić o kilka informacji
\VUgras{naczelnika stacji} \textsuperscript{(13)}, który rozmawiał z \VUgras{komisarzem}
\textsuperscript{(12)} nadzoru i z dyżurnym \VUgras{żandarmem} \textsuperscript{(11)}.

%%K t12-tit-1 sub
\VUpk{10.2pt}{40}{Zapisywanie bagaży.}
\VUsaut{3.0mm}

%%K t12-05-1 p
5. --- Kupiwszy gazetę w \VUgras{księgarni} \textsuperscript{(14)}, kazałem zważyć swoje
\VUgras{bagaże} \textsuperscript{(17)}, które \VUgras{tragarz} \textsuperscript{(16)}
właśnie przywiózł na \VUgras{wózku bagażowym} \textsuperscript{(19)}; \VUgras{urzędnik}
\textsuperscript{(22)} przydzielony do \VUgras{wagi} \textsuperscript{(21)} oznajmił mi,
że mój kufer waży trzydzieści pięć kilogramów; to czyniło \VUgras{nadwagę} pięciu
kilogramów, za którą jestem winien dopłatę, przy \VUgras{okienku} bagażowym
\textsuperscript{(23)}.

%%K t12-06-1 p
6. --- Ponieważ byłem za wcześnie, miałem wolny czas, żeby przypatrzeć się ruchowi
podróżnych na dworcu. Jedni składali albo odbierali swój drobny bagaż w
\VUgras{przechowalni bagażu} \textsuperscript{(25)}; inni sprawdzali \VUgras{rozkłady
jazdy} \textsuperscript{(26)}. Przybywający podróżni śpieszyli do \VUgras{wyjścia}
\textsuperscript{(32)} ze swoją \VUgras{walizką} \textsuperscript{(24)} w ręku. Kilku
czekało przy \VUgras{wydawaniu bagażu} \textsuperscript{(27)} i okazywało swój
\VUgras{kwit} \textsuperscript{(29)}, żeby odebrać swoje kufry, \VUgras{skrzynie}
\textsuperscript{(28)} albo \VUgras{kosze} \textsuperscript{(30)}.

%%K t12-07-1 p
7. --- \VUgras{Urzędnicy akcyzy} \textsuperscript{(33)} starannie badali paczki, żeby
stwierdzić, czy ma się coś do zgłoszenia.

%%K t12-08-1 p
8. --- Niektórzy podróżni udali się do \VUgras{bufetu} \textsuperscript{(34)}, gdzie
\VUgras{kelner} \textsuperscript{(35)} gorliwie ich obsługiwał. Muszą się śpieszyć, bo
\VUgras{pociąg} \textsuperscript{(36)}, który jest pośpieszny, nie zostanie długo na
dworcu.

%%K t12-09-1 p
9. --- Potem udałem się na peron odjazdowy i poszukałem bezpośredniego \VUgras{wagonu}
\textsuperscript{(37)}, żeby nie być zmuszonym zmieniać wagonu w czasie podróży.
Wsiadłem do wozu korytarzowego, który zawiera \VUgras{przedział}
\textsuperscript{(38)} dla palących i przedział zarezerwowany dla «pań samotnych».

%%K t12-tit-2 sub
\VUpk{10.2pt}{40}{Odjazd w nocy.}
\VUsaut{3.0mm}

%%K t12-10-1 p
10. --- Wszedłszy po \VUgras{stopniu} \textsuperscript{(40)}, zamknąwszy
\VUgras{drzwiczki} \textsuperscript{(39)}, zwinąwszy \VUgras{zasłonkę}
\textsuperscript{(41)} i umieściwszy swój drobny bagaż na \VUgras{siatce}
\textsuperscript{(43)}, usiadłem na \VUgras{ławce} \textsuperscript{(42)}. Widząc, jak
\VUgras{latarnik} \textsuperscript{(44)} zapala lampy, pomyślałem, że będę musiał
spędzić noc w wagonie, i przywołałem urzędnika przydzielonego do wynajmu
\VUgras{poduszek} \textsuperscript{(45)}, żeby poprosić o jedną z nich.

%%K t12-11-1 p
11. --- Konduktor pociągu przyszedł sprawdzić bilety. Zapytałem go, dlaczego jeszcze nie
odjechaliśmy, bo jest właściwa godzina. Odpowiedział mi, że \VUgras{kierownik pociągu}
\textsuperscript{(46)} jest zajęty ustawianiem rowerów w \VUgras{wagonie bagażowym}
\textsuperscript{(47)}. Ponadto nie zmieniono jeszcze wszystkich ogrzewaczy do nóg
\textsuperscript{(48)}.

%%K t12-12-1 p
12. --- Ale mieliśmy wkrótce odjechać, bo \VUgras{palacz} \textsuperscript{(50)}, na
swoim \VUgras{tendrze} \textsuperscript{(49)}, brał węgiel, żeby podsycić ogień
\VUgras{lokomotywy} \textsuperscript{(54)}, a \VUgras{maszynista}
\textsuperscript{(51)} czekał tylko na sygnał \VUgras{zastępcy naczelnika stacji}
\textsuperscript{(52)}, który był na \VUgras{peronie} \textsuperscript{(53)}.

%%K t12-13-1 p
13. --- \VUgras{Kocioł} \textsuperscript{(56)} lokomotywy głucho dudnił; \VUgras{komin}
\textsuperscript{(55)} zwracał potoki dymu zmieszanego z \VUgras{parą}
\textsuperscript{(60)}. Przenikliwy \VUgras{gwizd} \textsuperscript{(59)} zabrzmiał i
\VUgras{tłoki} \textsuperscript{(58)} ruszyły, popychając \VUgras{koła}
\textsuperscript{(57)} ciężkiej maszyny.

%%K t12-tit-3 sub
\VUsaut{3.0mm}
\VUpk{10.2pt}{40}{Odjazd!}
\VUsaut{3.0mm}

%%K t12-14-1 p
14. --- Prędkość pociągu biegnącego po \VUgras{szynach} \textsuperscript{(64)} rosła;
\VUgras{perony} \textsuperscript{(65)} znikły; minęliśmy \VUgras{ustępy}
\textsuperscript{(66)}, a także skład wagonów odstawiony na drugim \VUgras{torze}
\textsuperscript{(63)}: \VUgras{wagony sypialne} \textsuperscript{(67)}, \VUgras{wagon
restauracyjny} \textsuperscript{(68)}, \VUgras{wagon pocztowy} \textsuperscript{(69)}.

%%K t12-noto-1 noto
\VUnotes{143.2mm}{%
\textit{(*) Uwaga.} --- \textit{Rozumie się, że opis podróży jest według granicy
przedwojennej.}%
}

%%K t12-15-1 p
15. --- Potem \VUgras{stacja towarowa} \textsuperscript{(71)} przesunęła się przed
naszymi oczami. Pod halami urzędnicy ładowali \VUgras{towary} \textsuperscript{(72)}
wysyłane pociągiem zwykłym. \VUgras{Manewrowi} \textsuperscript{(76)} blisko
\VUgras{zbiornika wody} \textsuperscript{(73)} przestawiali \VUgras{wagony bydlęce}
\textsuperscript{(75)} i \VUgras{wagony platformy} \textsuperscript{(79)}, żeby utworzyć
\VUgras{pociąg towarowy} \textsuperscript{(80)}.

%%K t12-16-1 p
16. --- Przejechaliśmy przez ostatni \VUgras{rozjazd} \textsuperscript{(78)}, przy
którym stał zwrotniczy; minęliśmy \VUgras{tarczę sygnałową} \textsuperscript{(86)} i
dworzec znikł w oddali.

%%K t12-17-1 p
17. --- Wkrótce przejechaliśmy \VUgras{żelazny most} \textsuperscript{(74)}, który
przechodzi nad \VUgras{rzeką} \textsuperscript{(81)}. Przejechawszy ten most,
jechaliśmy wzdłuż rzeki, po której potężne \VUgras{holowniki} \textsuperscript{(82)}
ciągnęły ciężkie \VUgras{barki towarowe} \textsuperscript{(83)}. Widzieliśmy ponadto
\VUgras{statek pasażerski} \textsuperscript{(84)}, kiedy dobijał do \VUgras{pontonu}
\textsuperscript{(85)}.

%%K t12-18-1 p
18. --- Minąwszy z ogromną prędkością kilka stacji, przejechaliśmy przez kilka
\VUgras{tuneli} \textsuperscript{(87)} i zostawiliśmy za sobą niezliczone
\VUgras{domki} \textsuperscript{(88)} \VUgras{dróżników.} Ukołysany trzęsieniem
pociągu, zasnąłem głęboko.

%%K t12-tit-4 sub
\VUsaut{3.0mm}
\VUpk{10.2pt}{40}{Stacja Deutsch-Avricourt.}
\VUsaut{3.0mm}

%%K t12-19-1 p
19. --- Zostałem nagle obudzony głosem urzędnika, który wołał: «Deutsch-Avricourt!
\textsuperscript{(*)}. Wszyscy wysiadać!» Odbyła się inspekcja celna i wszystkie nudne
formalności graniczne. Musiałem uzgodnić swój zegarek kieszonkowy z \VUgras{wielkim
zegarem} \textsuperscript{(89)} dworca, który śpieszył się o pięćdziesiąt pięć minut;
ledwie obudzony, z trudem odróżniałem \VUgras{wielką wskazówkę} \textsuperscript{(90)}
od \VUgras{małej} \textsuperscript{(91)}; sama \VUgras{tarcza} \textsuperscript{(92)}
była zupełnie zamglona przez poranną mgiełkę.

%%K t12-20-1 p
20. --- Podczas gdy celnicy robili swoją inspekcję, ziewając odcyfrowywałem
\VUgras{afisze} \textsuperscript{(93)}, które zdobią mury dworca. Zjadłem śniadanie, żeby
zabić czas, sprawdziłem mapę niemieckiej \VUgras{sieci} \textsuperscript{(94)}, żeby
zobaczyć, ile odległości jeszcze dzieli mnie od Strasburga, i wróciłem do swojego
wagonu, wzmocniony nadzieją, że wkrótce osiągnę cel swojej podróży.
\VUsaut{6.0mm}
\VUfilet{22mm}
